\subsection{Lời gọi hệ thống -- System Call}

\textbf{\texttt{sys\_kilall}}
Chức năng: Đọc tên process từ \texttt{REGIONID}, so sánh với tên của các process trong multi-level queue và \texttt{running\_list}, giải phóng và terminate các process có cùng tên với tiến trình được truyển vào.

Trước tiên, ta đọc tên process từ \texttt{REGIONID} và in ra màn hình:

\begin{lstlisting}[style=Cstyle,language=C]
int __sys_killall(struct pcb_t *caller, struct sc_regs *regs)
{
    char proc_name[100];
    uint32_t data;
    uint32_t memrg = regs->a1;

    int i = 0;
    data = 0;
    while (data != -1) 
    {
        libread(caller, memrg, i, &data);
        proc_name[i] = data;
        if (data == -1)
            proc_name[i] = '\0';
        i++;
    }
    printf("The procname retrieved from memregionid %d is \"%s\"\n", memrg, proc_name);
\end{lstlisting}

\textbf{Giải thích:} Lấy \texttt{REGIONID} từ thanh ghi a1, đây là vùng nhớ chứa tiến trình cần kill. Duyệt từng byte trong vùng nhớ để đọc tên trình, lưu vào \verb|proc_name|.

Sau đó, ta lần lượt duyệt các multi-level queue:
\begin{lstlisting}[style=Cstyle,language=C]
    // MLQ
    for (int prio = 0; prio < MAX_PRIO; prio++)
    {
        struct queue_t *queue = &caller->mlq_ready_queue[prio];

        if (empty(queue)) continue;

        struct queue_t tmp_queue = {.size = 0};

        while (!empty(queue))
        {
            struct pcb_t *proc = dequeue(queue);
            char *pname = strrchr(proc->path, '/');
            pname = (pname != NULL) ? pname + 1 : proc->path;

            if (strcmp(pname, proc_name) != 0)
            {
                enqueue(&tmp_queue, proc);
            }
            else
            {
                free(proc);
            }
        }
        *queue = tmp_queue;
    }
\end{lstlisting}

\textbf{Giải thích:}

\begin{itemize}
    \item Duyệt từng queue trong nhóm các multi-level queue thông qua chỉ số ưu tiên
    \item Tạo 1 queue tạm là \verb|tmp_queue| để lưu các tiến trình ban đầu trong hàng đợi.
    \item Duyệt qua các tiến trình trong queue, với mỗi tiến trình, lấy tên của nó trong trường \verb|path|, lưu vào \verb|pname|.
    \item So sánh tên tiến trình \verb|pname| với \verb|proc_name|, nếu trùng thì gọi hàm \verb|free()| để kết thúc tiến trình, nếu không trùng thì giữ lại bằng cách enqueue vào tmp
    \item Gán lại queue gốc bằng \verb|tmp_queue|.
\end{itemize}

Ta xứ lí tương tự với \verb|running_list|
\begin{lstlisting}[style=Cstyle,language=C]
    // Running_list
    struct queue_t *runlist = caller->running_list;
    if (runlist != NULL && !empty(runlist))
    {
        struct queue_t tmp_runlist = {.size = 0};

        while (!empty(runlist))
        {
            struct pcb_t *proc = dequeue(runlist);
            char *pname = strrchr(proc->path, '/');
            pname = (pname != NULL) ? pname + 1 : proc->path;

            if (strcmp(pname, proc_name) != 0)
            {
                enqueue(&tmp_runlist, proc);
            }
            else
            {
                free(proc);
            }
        }

        *runlist = tmp_runlist;
    }

    return 0;
}
\end{lstlisting}